
\section{Le contexte philosophique}
\subsection{Ontologies : les principales approches}

Ontologie = Branche de la philosophie qui a pour objet l’être en tant qu’être
(incluant l’existence, la réalité) et qui étudie les propriétés
générales de l’être.
\subsubsection{Antiquité}
« Réalisme des idées » (Platon, « idéalistes médiévaux »), Nominalisme
$\neq$ Empirisme (Locke – c.f. I.2.c) [les sens, la matérialité]
\subsubsection{Réalisme}[défini aussi par Bitbol ; d’Espagnat ; Laloë]
(Pythagore, Démocrite, Newton, Galilée, Descartes, Diderot, Comte-Sponville)
Pythagorisme
(Néo-)Matérialisme

\subsubsection{Anti-réalisme}
réalité subjective (langage,...)
(Dummett, Berkeley)
\subsubsection{Idéalisme}
Kant, Hegel, Schopenhauer)
Versions de l’idéalisme...
modérée (Petitot)
radicale (Wittgenstein, Carnap, Quine, Putnam, Cassirer)
\subsubsection{Phénoménologie / Idéalisme transcendental(e)}
(Husserl, Heidegger)
Sur les fondements de la science, au-delà de la question des choses, des
phénomènes (kantiens). Rôle majeur de l’intuition, dans l’accession
envisageable au « noumène ».
\subsubsection{notes}
* L’observation et la prise de conscience créent le monde réel {\footnotesize Possible également, en toute logique, dans le réalisme au niveau du monde perçu (car la réalité est indépendante de
nous-même donc de notre perception ET de nos pensées).} – chez Fichte (réel créé par le
« moi »), Kant (nous contribuons activement par son exploration à la genèse du « monde »),
Hegel (objet créé par le sujet, le « moi ») et Schopenhauer (plus de séparation entre l’objet
et le sujet) $\Rightarrow$ vision circulaire de Hegel : « l’esprit, l’absolu » tend à se
 comprendre lui-même $\Rightarrow$ auto-déploiement (réalisation de soi) : savoir absolu
 au-delà d’une simple définition.
\subsection{Épistémologies : processus scientifiques}
{\footnotesize Épistémologies : au sens de la théorie des connaissances, de sa méthodologie.}

\begin{enumerate}[label=\Alph*]
  \item - Scepticisme [ni expérience (a posteriori), ni connaissance (a priori)] (Hume)
  \item - Phénomén(al)isme : existence d’objets extérieurs à nous mais réduits
à leurs apparences sensibles (pas de noumène) (Mach)
  \item - Empirisme pré et post-Kantien (Bacon, Locke)
versus rationalisme (Descartes, Leibniz)
  \item - Positivisme (Comte) : Sciences humaines $\Rightarrow$ Sciences exactes (= fondements)
  \item - Constructivisme (Goodman) : image/réalité biaisée par interaction esprit-réalité
  \item - Conventionnalisme (Poincaré) : libres (//exp.) décisions au coeur des sciences,
abstraction insaisissable dans l’activité de recherche
(hist. géométries non euclidiennes)
  \item - Descriptions initiales comme progrès continus de la Science,
et changements de ‘paradigmes’ (Bachelard, Kuhn)
\end{enumerate}
